\subsection{FRM Exam 2009: peakedness and left-tail statement}

\textbf{Enunciado}

\emph{An analyst gathered information about the return distributions for two portfolios during the same period.
The analyst states that Portfolio A is more peaked than a normal distribution and Portfolio B has a long tail on the left side. Which is correct?}


\[
\begin{array}{c|cc}
\text{Portafolio} & \text{Skewness} & \text{Kurtosis} \\
\hline
A & -1.6 & 1.9 \\
B & 0.8 & 3.2
\end{array}
\]


\begin{enumerate}[label=\alph*.]
\item The analyst’s assessment is correct.
\item The analyst’s assessment is correct for Portfolio A and incorrect for Portfolio B.
\item The analyst’s assessment is not correct for Portfolio A but is correct for Portfolio B.
\item \hl{The analyst’s assessment is incorrect for both portfolios.}
\end{enumerate}

\subsubsection*{Solución}

\begin{itemize}
\item ``More peaked than normal'' es leptocurtosis (kurtosis/exceso positivo).
\item ``Long left tail'' es skewness negativa.
\end{itemize}

\vspace{1cm}
Recordemos lo siguiente:
\[
\text{Skewness}<0 \text{ es cola izquierda más larga},
\qquad
\text{Skewness}>0 \text{ es cola derecha más larga}.
\]

Además, la \(\text{Kurtosis de una normal}=3.\) Por tanto:
\[
\text{Kurtosis}>3 \text{ es más picuda que la normal},
\qquad
\text{Kurtosis}<3 \text{ es menos picuda que la normal}.
\]

Para el portafolio \(A\):
\[
\text{Skewness}=-1.6,\qquad \text{Kurtosis}=1.9.
\]
El sesgo negativo indica una cola larga a la izquierda, pero como
\[
1.9<3,
\]
su distribución \emph{no} es más picuda que la normal. Entonces, la afirmación
del analista sobre \(A\) es falsa.

Para el portafolio \(B\):
\[
\text{Skewness}=0.8,\qquad \text{Kurtosis}=3.2.
\]
El sesgo positivo indica una cola larga a la derecha, no a la izquierda.
Entonces, la afirmación del analista sobre \(B\) también es falsa.




